\chapter{The results file}

\noindent This is \texttt{proofs/} verbatim, with one addition per environment: the name of the Lean
declaration (\verb|\lean|), the plain-label dependency list (\verb|\uses|), and \verb|\leanok| where
the item is complete --- on the definitions, which are Lean definitions, and on the results imported
from [BCIKS20], which are axioms. Every other statement is declared in \texttt{lean/} with a
\texttt{sorry} proof.

\section{Shared definitions}

\paragraph{Conventions.} These fix every symbol used below; the paper's \S3.1 is identical.
\begin{itemize}
\item $\N:=\{0,1,2,\dots\}$, $\Z$ the integers, $[n]:=\{1,\dots,n\}$,
$\frac1a\Z:=\{z/a:\ z\in\Z\}$ for an integer $a\ge1$, and $\max\emptyset:=-\infty$.
\item $\F$ is a finite field, $\F^*:=\F\setminus\{0\}$, and $\E$ a finite extension of $\F$.
\item For $0\le i\le m$, a multi-index at level $i$ is $\alpha=(\alpha_{i+1},\dots,\alpha_m)\in\N^{m-i}$,
$x^\alpha:=x_{i+1}^{\alpha_{i+1}}\cdots x_m^{\alpha_m}$, and for a weight vector $w\in\N^m$,
$\langle w,\alpha\rangle:=\sum_{l=i+1}^mw_l\alpha_l$. $e_l\in\N^m$ is the $l$-th unit vector and
$\mathbf 1:=(1,\dots,1)$.
\item For a commutative $\F$-algebra $R$, every $q\in R[x_{i+1},\dots,x_m]$ is uniquely
$\sum_\alpha q_\alpha x^\alpha$ with $q_\alpha\in R$ (finitely many nonzero), and
$\deg_{x_l}q:=\max\{\alpha_l:\ q_\alpha\ne0\}$.
\item $\F[T]_{\le e}:=\{g:\deg g\le e\}$ and $\F[T]_{<\lambda}:=\{g:\deg g<\lambda\}$, with $\deg0:=-\infty$.
\item A word on a finite set $L$ is a map $u:L\to A$; $P|_L$ is $x\mapsto P(x)$.
\item A weight function on $L$ is $\mu:L\to[0,1]$; $\mu(A):=\frac1{|L|}\sum_{x\in A}\mu(x)$ for
$A\subseteq L$; $\mu$ has common denominator $M\ge1$ if $\mu(x)\in\frac1M\Z$ for all $x$.
\end{itemize}

\begin{definition}[Field and domains]\label{def:domains}\lean{SymFRI.Tower}\leanok
$\F$ is a finite field; $k\ge 2$, $m\ge 1$ and $N$ are integers with $k^m\mid N$ and
$N\mid |\F|-1$. $L_0\subset\F^*$ is a coset of the subgroup of order $N$, and
$L_i:=\{x^k: x\in L_{i-1}\}$ for $1\le i\le m$, so $|L_i|=N/k^i$. The \emph{fibre} of
$y\in L_i$ is $\fib(y):=\{x\in L_{i-1}: x^k=y\}$; it equals $x_0\Omega_k$ for any
$x_0\in\fib(y)$, where $\Omega_k:=\{\zeta\in\F:\ \zeta^k=1\}$ is the group of $k$-th roots of unity, and $|\fib(y)|=k$.
\end{definition}

Since $k\mid N\mid|\F|-1$, the cyclic group $\F^*$ has an element $\omega$ of order $k$, so
$\Omega_k$ is as in \ref{lem:rootsum}; and $k\ne0$ in $\F$, because
$\operatorname{char}\F\nmid|\F|-1$. Consequently $k^i|L_i|=N$ for $0\le i\le m$, the rate of
$\RS[\F,L_i,k^{m-i}]$ is $\rho:=k^m/N=k^{m-i}/|L_i|$ for every $i$, and $|L_i|\ge k^{m-i}$.

\begin{definition}[Polynomials]\label{def:poly}\lean{SymFRI.InBox}\leanok
$R$ denotes a commutative $\F$-algebra (in practice $R=\F$ or $R=\F[T]$). For
$1\le i\le m+1$, $R[x_i,\dots,x_m]_{<k}:=\{q\in R[x_i,\dots,x_m]:\ \deg_{x_l}q<k\text{ for }i\le l\le m\}$.
\end{definition}

\begin{definition}[Kronecker map]\label{def:kron}\lean{SymFRI.kron}\leanok
For $0\le i\le m$ and $q\in R[x_{i+1},\dots,x_m]_{<k}$,
\[
U^{(i)}_q(Y):=q\big(Y,Y^{k},\dots,Y^{k^{m-i-1}}\big)\in R[Y],\qquad U_f:=U^{(0)}_f .
\]
For $i=m$, $q$ is a constant and $U^{(m)}_q=q$.
\end{definition}

\begin{definition}[Fold]\label{def:fold}\lean{SymFRI.fold}\leanok
For $u:L_{i-1}\to R$ and $y\in L_i$, the \emph{components} of $u$ are
\[
u_j(y):=\frac1k\sum_{x\in\fib(y)}x^{-j}\,u(x)\qquad(0\le j<k),
\]
and for $c\in R$ the \emph{fold} is $\Fold_c(u):L_i\to R$,
$\Fold_c(u)(y):=\sum_{j<k}c^j\,u_j(y)$.
\end{definition}

\begin{definition}[Weighted and graded degree]\label{def:wdeg}\lean{SymFRI.wdeg}\leanok
Let $w\in\N^m$. For $f=\sum_\alpha f_\alpha x^\alpha\in\F[x_1,\dots,x_m]$,
$\deg_w f:=\max\{\langle w,\alpha\rangle:\ f_\alpha\ne0\}$. For $0\le i\le m$ and
$q=\sum_{a,\alpha}q_{a,\alpha}T^ax^\alpha\in\F[T][x_{i+1},\dots,x_m]$ ($q_{a,\alpha}\in\F$),
$\gdeg_w q:=\max\{a+\langle w,\alpha\rangle:\ q_{a,\alpha}\ne0\}$.
\end{definition}

\begin{definition}[Codes]\label{def:codes}\lean{SymFRI.RS}\leanok
$\RS[\F,L,d]:=\{P|_L:\deg P<d\}$. The $t$-interleaved code $\RS[\F,L,d]^t$ is the set
of words $L\to\F^t$ whose $t$ coordinates lie in $\RS[\F,L,d]$; we identify
$\F[T]_{<t}$ with $\F^t$ through coefficients. The \emph{graded codes} are
\[
\Gc_i^{w,D}:=\big\{\,U^{(i)}_q\big|_{L_i}\ :\ q\in\F[T][x_{i+1},\dots,x_m]_{<k},\ \gdeg_w q\le D\,\big\},
\]
and the \emph{degree code} $C_{w,D}$ is the set of $T$-free elements of $\Gc_0^{w,D}$.
\end{definition}

\begin{definition}[Weighted agreement]\label{def:agree}\lean{SymFRI.agree}\leanok
For words $u,v$ on $L$ and $\mu:L\to[0,1]$,
$\agree_\mu(u,v):=\frac1{|L|}\sum_{x:\,u(x)=v(x)}\mu(x)$ (as in BCIKS20, \S7). For a set
$C$ of words, $\agree_\mu(u,C):=\max_{v\in C}\agree_\mu(u,v)$. Without a subscript,
$\mu\equiv1$: $\agree(u,v)=1-\Delta(u,v)$.
\end{definition}

Put $\Delta(u,v):=1-\agree(u,v)$; then $\Delta(u,w)\ge\Delta(v,w)-\Delta(v,u)$, because
$\{x:u(x)=w(x)\}\supseteq\{x:v(x)=w(x)\}\cap\{x:v(x)=u(x)\}$.

\paragraph{The caps.} For $w\in\N^m$ and $D\in\N$ put $d_0:=0$ and
$d_i:=\min\big(D,(k-1)\sum_{l\le i}w_l\big)$ for $1\le i\le m$; these are the formats of
\ref{def:protocol}.

\section{Tools}

\begin{lemma}[Vanishing on many points]\label{lem:roots}
Let $D\ge1$, let $P\in\F[Y]$ with $\deg P<D$, and let $A\subseteq\F$ with $|A|\ge D$. If
$P(y)=0$ for every $y\in A$, then $P=0$. Consequently two distinct $P,Q\in\F[Y]$ of degree $<D$
agree on fewer than $D$ points of $\F$.
\end{lemma}

\begin{proof}
A nonzero polynomial of degree $d$ over the field $\F$ has at most $d$ roots, and $d<D\le|A|$.
For the consequence apply this to $P-Q$.
\end{proof}

\begin{lemma}[Sum over the $k$-th roots of unity]\label{lem:rootsum}
Let $k\ge1$, let $\omega\in\F$ have multiplicative order exactly $k$, and let
$\Omega_k:=\{\omega^a:0\le a<k\}$. Then for every $j\in\Z$,
$\sum_{\zeta\in\Omega_k}\zeta^j=k$ if $k\mid j$, and $=0$ otherwise.
\end{lemma}

\begin{proof}
If $k\mid j$, every term is $1$. Otherwise $\omega^j\ne1$, and
$(\omega^j-1)\sum_{a<k}\omega^{aj}=\omega^{kj}-1=0$, so the sum is $0$.
\end{proof}

\begin{lemma}[Degree of a substitution]\label{lem:subst}\uses{def:wdeg}
Let $R$ be a commutative ring, $w\in\N^m$, $D\in\N$ and $I\subseteq\{1,\dots,m\}$. For $l\in I$
let $g_l\in R[T]$ with $\deg g_l\le w_l$. Let $q\in R[T][x_1,\dots,x_m]$ with $\gdeg_wq\le D$ and
$\deg_{x_l}q<k$ for every $l\notin I$, and let $\bar q$ be obtained by substituting $x_l\mapsto g_l$
for $l\in I$. Then $\gdeg_w\bar q\le D$ and $\deg_{x_l}\bar q<k$ for every $l\notin I$.
\end{lemma}

\begin{proof}
A monomial $T^ax^\alpha$ of $q$ maps to
$T^a\prod_{l\in I}g_l^{\alpha_l}\prod_{l\notin I}x_l^{\alpha_l}$. This is a sum of monomials of
$T$-degree at most $a+\sum_{l\in I}w_l\alpha_l$ and $x_l$-degree $\alpha_l$ for $l\notin I$, so
each has graded degree at most $a+\langle w,\alpha\rangle\le D$ and $x_l$-degree $<k$ for
$l\notin I$. Both bounds pass to the sum over the monomials of $q$.
\end{proof}

\section{The Kronecker map}

\begin{lemma}[Kronecker map]\label{lem:kron}\lean{SymFRI.kron_bijOn}\uses{def:kron}
For any commutative ring $R$, the map $q\mapsto U^{(i)}_q$ is an $R$-linear bijection
\[
R[x_{i+1},\dots,x_m]_{<k}\longrightarrow R[Y]_{<k^{m-i}},
\]
sending $x^\alpha$ to $Y^{n(\alpha)}$ with $n(\alpha):=\sum_{l>i}\alpha_l\,k^{\,l-i-1}$.
\end{lemma}

\begin{proof}
Substitution is $R$-linear and sends $x^\alpha$ to $Y^{n(\alpha)}$. Since $0\le\alpha_l<k$,
the integers $\alpha_{i+1},\dots,\alpha_m$ are the base-$k$ digits of $n(\alpha)$, so
$\alpha\mapsto n(\alpha)$ is a bijection onto $\{0,\dots,k^{m-i}-1\}$. A basis is sent to a
basis.
\end{proof}

\begin{lemma}[Evaluation is injective]\label{lem:inj}\lean{SymFRI.eval_injective}\uses{def:domains, lem:kron, lem:roots}
Let $R=\F[T]$ and $0\le i\le m$. The map $q\mapsto U^{(i)}_q\big|_{L_i}$ on
$R[x_{i+1},\dots,x_m]_{<k}$ is injective.
\end{lemma}

\begin{proof}
Write $U^{(i)}_q=\sum_a T^aP_a(Y)$ with $P_a\in\F[Y]_{<k^{m-i}}$ (\ref{lem:kron}). If
$U^{(i)}_q$ vanishes on $L_i$, then each $P_a$ vanishes on $L_i$. Since
$|L_i|=N/k^i\ge k^{m-i}>\deg P_a$, \ref{lem:roots} with $D=k^{m-i}$ gives $P_a=0$, so
$U^{(i)}_q=0$ and $q=0$ by \ref{lem:kron}.
\end{proof}

\section{Schwartz--Zippel in the total degree}

\begin{lemma}[Schwartz--Zippel for the total degree]\label{lem:sz-total}\lean{SymFRI.sz_total}\uses{def:wdeg, lem:roots}
Let $m\ge1$, $D\in\N$, and let $f\in\F[x_1,\dots,x_m]$ be nonzero with
$\deg_{\mathbf 1}f\le D$, where $\deg_{\mathbf 1}:=\deg_w$ for $w=\mathbf 1$ (\ref{def:wdeg}). For a
uniform $r\in\F^m$,
\[
\Pr\big[f(r)=0\big]\le\frac{D}{|\F|}.
\]
\end{lemma}

\begin{proof}\begin{enumerate}
\item Induction on $m$. For $m=1$, \ref{lem:roots} with $D=\deg f+1$ bounds the number of
roots of the nonzero $f$ by $\deg f\le D$, so the claim holds.
\item For $m\ge2$ write $f=\sum_{j\le d}f_j(x_1,\dots,x_{m-1})x_m^j$ with $f_d\ne0$ and $d\ge0$.
\item Then $\deg_{\mathbf 1}f_d+d\le\deg_{\mathbf 1}f$. Write $r=(r',r_m)$ with $r'$ uniform on $\F^{m-1}$; by the induction hypothesis applied to the $(m-1)$-variate polynomial $f_d$, $\Pr[f_d(r')=0]\le\deg_{\mathbf 1}f_d/|\F|\le(\deg_{\mathbf 1}f-d)/|\F|$.
\item Conditioned on $f_d(r')\ne0$ the polynomial $T\mapsto f(r',T)$ is nonzero of degree $d$ (for $d=0$ it is the nonzero constant $f_0(r')$), so it vanishes at the uniform $r_m$ with probability at most $d/|\F|$.
\item Adding the two bounds of steps 3 and 4 gives $\Pr[f(r)=0]\le(\deg_{\mathbf 1}f-d+d)/|\F|=\deg_{\mathbf 1}f/|\F|\le D/|\F|$.
\end{enumerate}\end{proof}

\begin{corollary}[The box bound]\label{cor:box}\lean{SymFRI.box_bound}\uses{lem:sz-total, def:wdeg, def:poly}
Let $0\ne f\in\F[x_1,\dots,x_m]$ with $\deg_{x_l}f<k$ for every $l$ (the box
$\F[x_1,\dots,x_m]_{<k}$ of \ref{def:poly}). Then $\deg_{\mathbf 1}f\le m(k-1)$, and for the
uniform $r\in\F^m$, $\Pr[f(r)=0]\le m(k-1)/|\F|$.
\end{corollary}

\begin{proof}
$\deg_{\mathbf 1}f=\max\{\sum_l\alpha_l:f_\alpha\ne0\}\le m(k-1)$; apply \ref{lem:sz-total} with
$D=m(k-1)$.
\end{proof}

\section{Fold equals substitution}

\begin{lemma}[Components of a polynomial]\label{lem:fold-parts}\lean{SymFRI.components_eq}\uses{def:fold, def:domains, lem:rootsum}
Let $1\le i\le m$, let $R$ be a commutative $\F$-algebra, and let
$P(Y)=\sum_{j<k}Y^jP_j(Y^k)$ with $P_j\in R[Y]$. Then the components of $u:=P|_{L_{i-1}}$
(\ref{def:fold}) are $u_j=P_j|_{L_i}$ for $0\le j<k$.
\end{lemma}

\begin{proof}
Fix $y\in L_i$ and $x_0\in\fib(y)$, so $\fib(y)=x_0\Omega_k$; then
\[
\sum_{x\in\fib(y)}x^{-j}P(x)=\sum_{j'<k}P_{j'}(y)\,x_0^{\,j'-j}\sum_{\zeta\in\Omega_k}
\zeta^{\,j'-j}=k\,P_j(y),
\]
the inner sum being $k$ for $j'=j$ and $0$ otherwise by \ref{lem:rootsum}, and $k$ is invertible by
\ref{def:domains}.
\end{proof}

\begin{lemma}[Fold equals substitution]\label{lem:fold-sub}\lean{SymFRI.fold_subst}\uses{lem:fold-parts, def:kron, lem:kron}
Let $1\le i\le m$, $q\in R[x_i,\dots,x_m]_{<k}$ and $u:=U^{(i-1)}_q|_{L_{i-1}}$. Then for every
$c\in R$,
\[
\Fold_c(u)=U^{(i)}_{q(c,\,x_{i+1},\dots,x_m)}\big|_{L_i}.
\]
\end{lemma}

\begin{proof}
Write $q=\sum_{j<k}x_i^jq_j$; by \ref{def:kron} $U^{(i-1)}_q(Y)=\sum_jY^jU^{(i)}_{q_j}(Y^k)$, so
$u_j=U^{(i)}_{q_j}|_{L_i}$ by \ref{lem:fold-parts}, and
$\Fold_c(u)=\sum_jc^jU^{(i)}_{q_j}|_{L_i}=U^{(i)}_{\sum_jc^jq_j}|_{L_i}$.
\end{proof}

\begin{lemma}[Inversion]\label{lem:fold-invert}\lean{SymFRI.fold_invert}\uses{def:fold, def:domains, lem:rootsum}
Let $1\le i\le m$ and let $u:L_{i-1}\to R$ have components $u_0,\dots,u_{k-1}$
(\ref{def:fold}). Then $u(x)=\sum_{j<k}x^j\,u_j(x^k)$ for every $x\in L_{i-1}$.
\end{lemma}

\begin{proof}
Put $y:=x^k$; then $\sum_{j<k}x^ju_j(y)=\frac1k\sum_{x'\in\fib(y)}u(x')\sum_{j<k}(x/x')^j=u(x)$,
since $x/x'\in\Omega_k$ and the inner sum is $k$ for $x'=x$ and $0$ otherwise
(\ref{lem:rootsum}).
\end{proof}

\begin{corollary}[Honest terminal]\label{cor:terminal}\lean{SymFRI.terminal}\uses{lem:fold-sub}
Let $R=\F[T]$, $f\in\F[x_1,\dots,x_m]_{<k}$, $w\in\N^m$, $r\in\F^m$ and
$c_i:=r_iT^{w_i}$. Then $\Fold_{c_m}\circ\cdots\circ\Fold_{c_1}\big(U_f|_{L_0}\big)$ is the
constant $g(T):=f\big(r_1T^{w_1},\dots,r_mT^{w_m}\big)$.
\end{corollary}

\begin{proof}
Induction on $i$ with \ref{lem:fold-sub}; at $i=m$, $U^{(m)}$ is the identity.
\end{proof}

\section{Folding preserves the graded codes}

\begin{lemma}[Stability]\label{lem:stab}\lean{SymFRI.stab}\uses{lem:fold-sub, lem:subst, def:codes}
Let $1\le i\le m$, $w\in\N^m$, $D\in\N$, $r\in\F$ and $c:=rT^{w_i}$. Then
$\Fold_c\big(\Gc_{i-1}^{w,D}\big)\subseteq\Gc_i^{w,D}$.
\end{lemma}

\begin{proof}
Let $v=U^{(i-1)}_q|_{L_{i-1}}\in\Gc_{i-1}^{w,D}$ with $\gdeg_wq\le D$. By \ref{lem:fold-sub},
$\Fold_c(v)=U^{(i)}_{q(c,x_{i+1},\dots,x_m)}|_{L_i}$, and $q(c,x_{i+1},\dots,x_m)$ lies in the box
with graded degree $\le D$ by \ref{lem:subst}.
\end{proof}

\begin{lemma}[The honest $T$-degree]\label{lem:stab-cap}\lean{SymFRI.stab_cap}\uses{lem:fold-sub, lem:subst, def:codes, def:wdeg}
Let $1\le i\le m$, $w\in\N^m$, $D\in\N$, $c_l:=r_lT^{w_l}$ for $1\le l\le i$, and let
$f\in\F[x_1,\dots,x_m]_{<k}$ have $\gdeg_wf\le D$. Then
$\Fold_{c_i}\circ\cdots\circ\Fold_{c_1}\big(U_f|_{L_0}\big)$ is a word of $\Gc_i^{w,D}$ whose values
have $T$-degree at most $d_i$.
\end{lemma}

\begin{proof}
By \ref{lem:fold-sub} applied $i$ times the word is $U^{(i)}_q$ with
$q=f(c_1,\dots,c_i,x_{i+1},\dots,x_m)$, which lies in $\Gc_i^{w,D}$ by \ref{lem:subst}. A monomial
$x^\alpha$ of $f$ contributes the $T$-exponent $\sum_{l\le i}w_l\alpha_l\le(k-1)\sum_{l\le i}w_l$,
and this exponent is at most the graded degree, hence $\le D$.
\end{proof}

\section{Divisibility}

\begin{lemma}[Divisibility]\label{lem:div}\lean{SymFRI.divisible}\uses{def:codes, lem:inj, lem:roots}
Let $1\le i\le m$, $\gamma\in\N$, $v\in\Gc_i^{w,D}$ and $A\subseteq L_i$ with $|A|\ge k^{m-i}$.
If $T^\gamma$ divides $v(y)$ for every $y\in A$, then $v=T^\gamma p$ on $L_i$ for some
$p\in\Gc_i^{w,D-\gamma}$; for $\gamma>D$, $v=0$.
\end{lemma}

\begin{proof}
Write $v=U^{(i)}_q|_{L_i}$ with $\gdeg_w q\le D$ and $q=\sum_aT^aq_a$. For $a<\gamma$, the
$T^a$-coefficient of $v$ is $U^{(i)}_{q_a}|_{L_i}$, a codeword of $\RS[\F,L_i,k^{m-i}]$ that
vanishes on $A$; since $|A|\ge k^{m-i}>\deg U^{(i)}_{q_a}$, \ref{lem:roots} with $D=k^{m-i}$
gives $U^{(i)}_{q_a}|_{L_i}=0$, so $q_a=0$ by \ref{lem:inj}. Hence $q=T^\gamma Q$ with $\gdeg_w Q\le D-\gamma$, and $v=T^\gamma p$ for
$p:=U^{(i)}_Q|_{L_i}\in\Gc_i^{w,D-\gamma}$. If $\gamma>D$ then $q=0$.
\end{proof}

\section{Unfolding}

\begin{lemma}[Unfolding]\label{lem:unfold}\lean{SymFRI.unfold}\uses{def:kron, def:codes, lem:div, lem:inj, lem:fold-parts}
Let $1\le i\le m$, $w\in\N^m$, $D\in\N$. Let $v_0,\dots,v_{k-1}\in\Gc_i^{w,D}$ and
$p_0,\dots,p_{k-1}:L_i\to\F[T]$ satisfy $v_j(y)=T^{w_ij}p_j(y)$ for all $y\in L_i$ and
$0\le j<k$. Define $h:L_{i-1}\to\F[T]$ by $h(x):=\sum_{j<k}x^jp_j(x^k)$. Then
$h\in\Gc_{i-1}^{w,D}$, and the components of $h$ (\ref{def:fold}) are $h_j=p_j$.
\end{lemma}

\begin{proof}
Write $v_j=T^{w_ij}p_j$ on $L_i$; by \ref{lem:div} with $A=L_i$, \ref{lem:inj} gives
$v_j=U^{(i)}_{q'_j}|_{L_i}$ with $p_j=U^{(i)}_{q'_j}|_{L_i}$ and $\gdeg_w q'_j\le D-w_ij$.
Put $Q:=\sum_{j<k}x_i^{\,j}q'_j$. A monomial $x_i^{\,j}T^ax^\beta$ of $Q$ has graded degree
$w_ij+a+\sum_{l>i}w_l\beta_l\le w_ij+(D-w_ij)=D$, so $Q$ has graded degree $\le D$, and by
\ref{def:kron} $h=U^{(i-1)}_Q|_{L_{i-1}}\in\Gc_{i-1}^{w,D}$. Its components are $p_j$ by
\ref{lem:fold-parts}.
\end{proof}

\section{Interleaved codes as codes over an extension}

\begin{lemma}[Interleaved code $=$ extension-field code]\label{lem:ext}\lean{SymFRI.ext_linear}\uses{def:codes}
Let $\E\supseteq\F$ be a finite extension, $e_0,\dots,e_{[\E:\F]-1}$ an $\F$-basis of
$\E$, and $\varphi:\F^{[\E:\F]}\to\E$, $a\mapsto\sum_s a_se_s$, applied pointwise to words on
$L\subseteq\F$. Then $\varphi$ is an $\F$-linear bijection $\RS[\F,L,d]^{[\E:\F]}\to\RS[\E,L,d]$, and
$\varphi\big(\sum_{j\le l} z_ju_j\big)=\sum_{j\le l} z_j\,\varphi(u_j)$ for $z_0,\dots,z_l\in\F$.
\end{lemma}

\begin{proof}
Every $P\in\E[X]_{<d}$ decomposes uniquely as $P=\sum_s e_sP_s$ with $P_s\in\F[X]_{<d}$ (expand
each coefficient in the basis), and for $x\in L\subseteq\F$,
$P(x)=\sum_s e_sP_s(x)=\varphi\big((P_s(x))_s\big)$: this is the bijection, and
$\F$-linearity gives the identity for $z_j\in\F$.
\end{proof}

\begin{lemma}[$\varphi$ preserves agreement]\label{lem:ext-agree}\lean{SymFRI.ext_agree}\uses{def:codes, def:agree}
Let $\varphi:\F^t\to\E$ be an injective $\F$-linear map, applied pointwise to words on a set $L$.
Then $\{x:u(x)=v(x)\}=\{x:\varphi(u(x))=\varphi(v(x))\}$ for all words $u,v$; hence $\varphi$
preserves $\Delta$ and $\agree_\mu$.
\end{lemma}

\begin{proof}
$\varphi$ is injective, so $u(x)=v(x)\iff\varphi(u(x))=\varphi(v(x))$; the agreement set is
unchanged, and $\Delta,\agree_\mu$ are functions of it.
\end{proof}

\begin{remark}
The identity $\varphi\big(\sum_jz_ju_j\big)=\sum_jz_j\varphi(u_j)$ for $z_j\in\F$ is the only place
where it matters that the challenges lie in $\F$ and not in $\E$; it is what \ref{cor:share} uses.
\end{remark}

\section{The subcode lemma}

\begin{lemma}[Subcode lemma]\label{lem:subcode}\lean{SymFRI.subcode}
Let $\E\supseteq\F$, let $W$ be an $\F$-linear space of words $L\to\E$ and let $v_0,\dots,v_l$ be
words $L\to\E$. Let $S'\subseteq\F$ with $|S'|>l$. If $\sum_{j\le l}z^jv_j\in W$ for every
$z\in S'$, then $v_j\in W$ for every $j$.
\end{lemma}

\begin{proof}
Pick distinct $z_0,\dots,z_l\in S'$ and put $P_a:=\sum_jz_a^jv_j\in W$. The Vandermonde
matrix $\Lambda:=(z_a^{\,j})_{a,j}\in\F^{(l+1)\times(l+1)}$ is invertible, so
$v_j=\sum_a(\Lambda^{-1})_{j,a}P_a$ with $(\Lambda^{-1})_{j,a}\in\F$. Since $W$ is $\F$-linear,
$v_j\in W$.
\end{proof}

\section{Granularity}

\begin{lemma}[Granularity]\label{lem:gran}\lean{SymFRI.granularity}\uses{def:agree, def:domains}
Let $\mu:L_i\to\frac1{k^i}\Z$. Then $\agree_\mu(u,v)\in\frac1N\Z$ for all words
$u,v$ on $L_i$. Hence, for $\beta\in\mathbb R$ and
$\beta^+:=\min\{x\in\frac1N\Z:\ x>\beta\}$, one has $\agree_\mu(u,C)>\beta$ if and only if
$\agree_\mu(u,C)\ge\beta^+$, for every set $C$ of words on $L_i$.
\end{lemma}

\begin{proof}
Each $\mu(x)$ equals $m_x/k^i$ with $m_x\in\Z$, so
\[
\agree_\mu(u,v)=\frac1{|L_i|}\sum_{x:\,u(x)=v(x)}\frac{m_x}{k^i}=\frac1N\sum_{x:\,u(x)=v(x)}m_x,
\]
using \ref{def:domains}. The value lies in $\frac1N\Z$, and the equivalence is the definition of
$\beta^+$ on that lattice.
\end{proof}

\section{Fibre averages and lifting}

\begin{definition}[Fibre average]\label{def:fibre-avg}\lean{SymFRI.fibreAvg}\leanok\uses{def:domains}
Let $\mu$ be a weight on $L_{i-1}$. Its \emph{fibre average} is the weight $\nu$ on $L_i$ with
\[
\nu(y):=\frac1k\sum_{x\in\fib(y)}\mu(x)\qquad(y\in L_i).
\]
\end{definition}

\begin{lemma}[Lifting]\label{lem:lift}\lean{SymFRI.lift}\uses{def:agree, def:fibre-avg, def:domains}
Let $\mu$ be a weight on $L_{i-1}$ with fibre average $\nu$ (\ref{def:fibre-avg}). If $u,h$, maps
on $L_{i-1}$ into any set, agree on $\fib(y)$ for every $y\in A\subseteq L_i$, then
\[
\agree_\mu(u,h)\ \ge\ \nu(A):=\frac1{|L_i|}\sum_{y\in A}\nu(y).
\]
\end{lemma}

\begin{proof}
The set of $x\in L_{i-1}$ on which $u$ and $h$ agree contains the union of the fibres $\fib(y)$,
$y\in A$, and these are disjoint. Hence
$\agree_\mu(u,h)\ge\frac1{|L_{i-1}|}\sum_{y\in A}\sum_{x\in\fib(y)}\mu(x)$, and
$\frac1{|L_{i-1}|}\sum_{x\in\fib(y)}\mu(x)=\frac1{k|L_i|}\sum_{x\in\fib(y)}\mu(x)=\frac{\nu(y)}{|L_i|}$.
\end{proof}

\section{Counting bad challenge sequences}

\begin{lemma}[Counting bad challenge sequences]\label{lem:bad-count}\lean{SymFRI.bad_count}
For $1\le i\le m$ let $S_i:\F^{i-1}\to2^{\F}$ with $|S_i(r_1,\dots,r_{i-1})|\le\Theta_i$ for
every $(r_1,\dots,r_{i-1})$. Then
\[
\big|\{r\in\F^m:\ \exists i,\ r_i\in S_i(r_1,\dots,r_{i-1})\}\big|
\le|\F|^{m-1}\sum_{i=1}^m\Theta_i .
\]
\end{lemma}

\begin{proof}
For fixed $i$, the tuples with $r_i\in S_i(r_1,\dots,r_{i-1})$ number at most
$|\F|^{i-1}\cdot\Theta_i\cdot|\F|^{m-i}$; sum over $i$.
\end{proof}

\section{Imported results}

All results below are from Ben-Sasson, Carmon, Ishai, Kopparty, Saraf,
\emph{Proximity Gaps for Reed--Solomon Codes}, ePrint 2020/654 (version of 3 July 2021),
cited as [BCIKS20]. [BCIKS20] index Reed--Solomon codes by the degree bound, which is one less than the dimension,
and write the evaluation domain as a set of size $n$. Their code of dimension $d$ on the domain
$L$ is our $\RS[\F,L,d]$, with $n=|L|$ and the same rate $\rho=d/|L|$.

\begin{imported}[{[BCIKS20, Theorem 7.2]}: weighted correlated agreement over curves]\label{ext:bciks72}\lean{SymFRI.bciks72}\leanok
Let $V=\RS[\F,L,d]$, $n=|L|$, $\rho=d/n$. Let $u_0,\dots,u_l:L\to\F$, and let
$\mu:L\to[0,1]$ have common denominator $M$. Let $\kappa\ge3$ be an integer and
$\alpha\ge\alpha_0(\rho,\kappa):=\sqrt\rho\,(1+\tfrac1{2\kappa})$. Let
\[
S:=\Big\{z\in\F:\ \agree_\mu\Big(\textstyle\sum_{j\le l}z^ju_j,\,V\Big)\ge\alpha\Big\}
\]
and suppose
\[
|S|>\max\Big(\frac{(1+\frac1{2\kappa})^7\kappa^7}{3\rho^{3/2}}\,n^2\,l,\ \
\frac{2\kappa+1}{\sqrt\rho}\,(Mn+1)\,l\Big).\tag{$\star$}
\]
Then there are $v_0,\dots,v_l\in V$ with
$\mu\big(\{x\in L:\ u_j(x)=v_j(x)\ \forall j\}\big)\ge\alpha$.
\end{imported}

\begin{remark}[Two facts from the proof of \ref{ext:bciks72}]\label{ext:proof-facts}
\leavevmode
\begin{enumerate}
\item[(F1)] \emph{The witnesses are arbitrary.} The proof starts from an arbitrary choice,
for each $z\in S$, of a codeword $P_z\in V$ with
$\agree_\mu(\sum_jz^ju_j,P_z)\ge\alpha$ ([BCIKS20, p.~33, opening of \S7.1]; the
unweighted proofs make the same choice, pp.~15 and 19). The codewords $v_j$ produced
satisfy
\[
\textstyle\sum_{j}z^jv_j=P_z\qquad\text{for all }z\in S',
\]
for a subset $S'\subseteq S$ with $|S'|\ge(Mn+1)l$ (Proposition 5.5, eq.~(5.10), and
the paragraph on p.~34), for \emph{those same} $P_z$.
\item[(F2)] \emph{Only $|S|$ matters.} The proofs of Theorems 5.1, 6.2, 7.2,
Proposition 5.5 and Lemma A.1 use $S$ only through its cardinality and through
membership $z\in S$. The size of the field enters only (i) through the bound
$|S|\le|\F|$, used on p.~22 to find a point $x_0\in\F$ outside fewer than
$(\kappa+\frac12)^2n$ bad values, and (ii) through the separability hypothesis (char
$\F>D_Y$, where $D_Y<\frac{2\kappa+1}{2\sqrt\rho}$ bounds the $Y$-degree of the interpolating
polynomial $Q(X,Y,Z)$ of \S5, or Appendix~C) of \S5.2.3. No step requires $S$ to range over the whole field.
\end{enumerate}
\end{remark}

\begin{imported}[{[BCIKS20, Theorem 7.1]}: weighted correlated agreement over curves, version I]\label{ext:bciks71}\lean{SymFRI.bciks71}\leanok
Let $V=\RS[\F,L,d]$, $n=|L|$, $\rho=d/n$, $u_0,\dots,u_l:L\to\F$, and $\mu:L\to[0,1]$
with common denominator $M$. Let $\alpha\in(\sqrt\rho,1)$ and
$\eta:=\min(\alpha-\sqrt\rho,\frac{\sqrt\rho}{20})$. With $S$ as in \ref{ext:bciks72},
suppose
\[
|S|>l\,\varepsilon\,|\F|\quad\text{and}\quad|S|\ge l\,(Mn+1)\Big(\frac1\eta+\frac3{\sqrt\rho}\Big),\tag{$\star\star$}
\]
where $\varepsilon$ is the error of [BCIKS20, Theorem~1.2] at $\delta=1-\alpha$; in the
unique-decoding range $\alpha>\frac{1+\rho}2$ it is $\varepsilon=n/|\F|$, so the first
condition reads $|S|>l\,n$. Then there are $v_0,\dots,v_l\in V$ with
$\mu(\{x:\ u_j(x)=v_j(x)\ \forall j\})\ge\alpha$.
\end{imported}

\begin{remark}[Facts F1--F2 for \ref{ext:bciks71}]\label{ext:facts71}
The proof of \ref{ext:bciks71} is the same as that of \ref{ext:bciks72} (\S7.1 of
[BCIKS20]): it applies Theorems~6.1/6.2 to $S$ and Lemmas~7.5--7.6 to the identity
$\sum_jz^jv_j=P_z$ on $S'$. In the unique-decoding range, $S'=S$ (p.~34: ``$P_z$ and
$\sum_jz^jv_j$ are both decodings of $\sum_jz^ju_j$, and the decoding is unique'').
Hence (F1) and (F2) hold for \ref{ext:bciks71} as well.
\end{remark}

\section{The hypothesis}

\begin{hypothesis}\label{hyp:ca}\lean{SymFRI.CA}\uses{ext:bciks72, ext:proof-facts}
For $\alpha_*\in(0,1]$ and a function $\Theta(n,M,l)$, $\mathrm H(\alpha_*,\Theta)$ is the
following assertion. For every finite field $\E\supseteq\F$, every $L\subseteq\F$ with
$|L|=n$, every $V=\RS[\E,L,d]$ of rate $\rho=d/n$, all words $u_0,\dots,u_l:L\to\E$, every
weight $\mu$ whose values have denominator $M$, every $\alpha\ge\alpha_*$, and every choice of
witnesses $P_z\in V$ with $\agree_\mu\big(\sum_{j\le l}z^ju_j,P_z\big)\ge\alpha$ for the
elements $z$ of a set $S\subseteq\F$: if $|S|>\Theta(n,M,l)$, then there are $v_0,\dots,v_l\in V$
with
\[
\mu\big(\{x\in L:\ u_j(x)=v_j(x)\ \forall j\}\big)\ge\alpha
\quad\text{and}\quad
\textstyle\sum_{j\le l}z^jv_j=P_z\ \text{ for more than }l\text{ values }z\in S .
\]
\end{hypothesis}

{\renewcommand{\theproposition}{P-H1}
\begin{proposition}[H holds at the Johnson radius]\label{prop:H1}\lean{SymFRI.H_johnson}\uses{ext:bciks72, ext:proof-facts}
Let $\kappa\ge3$ and $\alpha_0:=\sqrt\rho\,(1+\frac1{2\kappa})$. Then
$\mathrm H(\alpha_0,\Theta_J)$ holds for
\[
\Theta_J(n,M,l):=l\,\max\Big(\frac{(1+\frac1{2\kappa})^7\kappa^7}{3\rho^{3/2}}\,n^2,\ \
\frac{2\kappa+1}{\sqrt\rho}\,(Mn+1)\Big).
\]
\end{proposition}

\begin{proof}
Let $\E,L,d,V,u_j,\mu,M,\alpha\ge\alpha_0,S,P_z$ be as in \ref{hyp:ca}, with
$|S|>\Theta_J(n,M,l)$. The proof of \ref{ext:bciks72} uses the set $S$ only through $|S|$ and
membership in it, and both places where the field enters hold over $\E$, since
$|S|\le|\F|\le|\E|$ and $\E$ has the characteristic of $\F$; so its conclusion holds for the
field $\E$ and the set $S$, giving $v_0,\dots,v_l\in V$ with the stated agreement. By the first
fact of \ref{ext:proof-facts} these $v_j$ satisfy $\sum_jz^jv_j=P_z$ for the given witnesses on
a set of more than $l$ values of $z$.
\end{proof}

{\renewcommand{\theproposition}{P-H2}
\begin{proposition}[H holds in the unique-decoding regime]\label{prop:H2}\lean{SymFRI.H_unique}\uses{ext:bciks71, ext:facts71}
Put $\eta_{\mathrm U}:=\min\big(\frac{(1-\sqrt\rho)^2}2,\frac{\sqrt\rho}{20}\big)$. For every
$\alpha_*>\frac{1+\rho}2$, $\mathrm H(\alpha_*,\Theta_U)$ holds for
\[
\Theta_U(n,M,l):=l\,(Mn+1)\Big(\frac1{\eta_{\mathrm U}}+\frac3{\sqrt\rho}\Big).
\]
\end{proposition}

\begin{proof}
Let $\alpha\ge\alpha_*>\frac{1+\rho}2$ and $|S|>\Theta_U(n,M,l)$, and put
$\eta:=\min(\alpha-\sqrt\rho,\frac{\sqrt\rho}{20})$. Since
$\alpha-\sqrt\rho>\frac{1+\rho}2-\sqrt\rho=\frac{(1-\sqrt\rho)^2}2$ we have
$\eta\ge\eta_{\mathrm U}$, so $|S|>l(Mn+1)(\frac1\eta+\frac3{\sqrt\rho})$; also
$\alpha>\frac{1+\rho}2$ gives $\varepsilon=n/|\F|$ in $(\star\star)$, and
$(Mn+1)(\frac1\eta+\frac3{\sqrt\rho})>n$. So \ref{ext:bciks71} applies, with the same
transposition to $\E$ as in \ref{prop:H1}, and \ref{ext:facts71} supplies the curve identity
for the chosen witnesses on more than $l$ values of $z$.
\end{proof}

\section{Correlated agreement for linear subcodes}

\begin{consequence}[Correlated agreement for subcodes]\label{cor:share}\lean{SymFRI.share}\uses{hyp:ca, lem:subcode}
Assume $\mathrm H(\alpha_*,\Theta)$, let $\E,L,d,V,u_0,\dots,u_l$ and $\mu,M$ be as in
\ref{hyp:ca}, and let $\alpha\ge\alpha_*$. Let $W\subseteq V$ be an $\F$-linear subspace and let
\[
S:=\Big\{z\in\F:\ \exists P_z\in W\ \text{with}\ \agree_\mu\Big(\textstyle\sum_{j\le l}z^ju_j,P_z\Big)\ge\alpha\Big\}
\]
satisfy $|S|>\Theta(n,M,l)$. Then there are $v_0,\dots,v_l\in W$ with
\[
\mu\big(\{x\in L:\ u_j(x)=v_j(x)\ \forall j\}\big)\ge\alpha .
\]
\end{consequence}

\begin{proof}
Fix $P_z\in W$ for $z\in S$ as in the definition of $S$, and apply \ref{hyp:ca}: it gives
$v_0,\dots,v_l\in V$ with the stated agreement and $\sum_jz^jv_j=P_z$ for more than $l$ values
$z\in S$. Those sums lie in $W$, so $v_j\in W$ by \ref{lem:subcode}.
\end{proof}

\section{The protocol and the query weights}

\begin{definition}[Degree test $\Pi_{w,D}$]\label{def:protocol}\lean{SymFRI.Prover}\leanok\uses{def:domains, def:kron, def:fold}
Parameters: $\F,k,m,N,L_0,\dots,L_m$ as in \ref{def:domains}; $w\in\N^m$; $D\in\N$;
$s\ge1$.
\begin{enumerate}
\item The prover sends $u_0:L_0\to\F$ (honestly $u_0=U_f|_{L_0}$).
\item For $i=1,\dots,m$: the verifier sends a uniform $r_i\in\F$; put $c_i:=r_iT^{w_i}$.
The prover sends $u_i:L_i\to\F[T]_{\le d_i}$ (honestly $u_i=\Fold_{c_i}(u_{i-1})$).
For $i=m$ the word $u_m$ is constant, sent as one polynomial $g\in\F[T]_{\le D}$.
\item Repeat $s$ times: pick a uniform $\xi_0\in L_0$, put $\xi_i:=\xi_{i-1}^k$, and check
$u_i(\xi_i)=\Fold_{c_i}(u_{i-1})(\xi_i)$ for $i=1,\dots,m$ (reading $u_{i-1}$ on
$\fib(\xi_i)$).
\end{enumerate}
The verifier accepts iff all checks pass. (The format $u_i(y)\in\F[T]_{\le d_i}$ and
$\deg g\le D$ is part of the message space.)
\end{definition}

\begin{definition}[Query weights, after {[BCIKS20, \S8.2.1]}]\label{def:weights}\lean{SymFRI.nu}\leanok\uses{def:fibre-avg, def:protocol, def:fold}
For a transcript, define $\mu^{(0)}\equiv1$ on $L_0$ and, for $1\le i\le m$, let $\nu^{(i)}$ be
the fibre average (\ref{def:fibre-avg}) of $\mu^{(i-1)}$ and put, for $y\in L_i$,
\[
\mu^{(i)}(y):=\begin{cases}\nu^{(i)}(y)&\text{if }u_i(y)=\Fold_{c_i}(u_{i-1})(y),\\0&\text{otherwise.}\end{cases}
\]
The values of $\mu^{(i)},\nu^{(i)}$ are multiples of $1/k^i$; hence, by \ref{def:domains},
every $\mu^{(i)}$- or $\nu^{(i)}$-agreement on $L_i$ is a multiple of $1/N$.
\end{definition}

\section{Acceptance as a weight}

\begin{imported}[{[BCIKS20, Claim 8.5]}: acceptance as a weight]\label{ext:claim85}\lean{SymFRI.claim85}\leanok
With the weights of \ref{def:weights}, the probability that one query of
\ref{def:protocol} accepts equals $\frac1{|L_m|}\sum_{y\in L_m}\mu^{(m)}(y)$.
\end{imported}
\noindent\emph{Note.} [BCIKS20] sample $y\in L_m$ uniformly and then a uniform element of
each fibre downwards; sampling a uniform $\xi_0\in L_0$ and setting $\xi_i:=\xi_{i-1}^k$ gives
the same distribution of query paths. Their proof (induction on $i$, p.~43) applies
verbatim.

\section{Starting the induction}

\begin{lemma}[Start]\label{lem:start}\lean{SymFRI.start}\uses{def:codes, def:agree, def:kron}
For every $u_0:L_0\to\F$, $\agree\big(u_0,\Gc_0^{w,D}\big)=\agree\big(u_0,C_{w,D}\big)$.
\end{lemma}

\begin{proof}
Both sides are maxima over words, and $C_{w,D}\subseteq\Gc_0^{w,D}$, so it suffices to show
that every $v\in\Gc_0^{w,D}$ has an element of $C_{w,D}$ agreeing with it wherever $u_0$ does.
Let $v=U^{(0)}_q|_{L_0}$ and let $q_0$ be the $T^0$-coefficient of $q$; then
$U^{(0)}_{q_0}|_{L_0}\in C_{w,D}$ and $U^{(0)}_q=U^{(0)}_{q_0}$ wherever the value is
$T$-free, in particular wherever it equals $u_0(x)\in\F$.
\end{proof}

\section{Agreement does not increase}

\begin{lemma}[No increase]\label{lem:no-increase}\lean{SymFRI.no_increase}\uses{def:weights, def:agree}
Let $1\le i\le m$ and let $\mu^{(i)},\nu^{(i)}$ be the weights of \ref{def:weights} for a
transcript. For every set $C$ of words on $L_i$,
\[
\agree_{\mu^{(i)}}\big(u_i,C\big)\ \le\ \agree_{\nu^{(i)}}\big(\Fold_{c_i}(u_{i-1}),C\big).
\]
\end{lemma}

\begin{proof}
By \ref{def:weights}, $\mu^{(i)}(y)=\nu^{(i)}(y)$ where
$u_i(y)=\Fold_{c_i}(u_{i-1})(y)$ and $\mu^{(i)}(y)=0$ elsewhere. Hence for every word $w$ on
$L_i$ the set $\{y:u_i(y)=w\}$ meets $\{y:\mu^{(i)}(y)>0\}$ inside
$\{y:\Fold_{c_i}(u_{i-1})(y)=w\}$, where $\mu^{(i)}=\nu^{(i)}$; so
$\agree_{\mu^{(i)}}(u_i,w)\le\agree_{\nu^{(i)}}(\Fold_{c_i}(u_{i-1}),w)$. Taking the maximum
over $w\in C$ gives the claim.
\end{proof}

\section{The round lemma}

\begin{lemma}[Descent]\label{lem:round-up}\lean{SymFRI.round_up}\uses{def:domains, def:fibre-avg, lem:div, lem:fold-invert, lem:lift, lem:unfold}
Let $1\le i\le m$, $w\in\N^m$, $D\in\N$, let $u:L_{i-1}\to\F[T]$ have components
$u_0,\dots,u_{k-1}$ (\ref{def:fold}), let $\mu$ be a weight on $L_{i-1}$ with fibre average $\nu$
(\ref{def:fibre-avg}), and let $A\subseteq L_i$ with $\nu(A)>\rho$. If there are
$v_0,\dots,v_{k-1}\in\Gc_i^{w,D}$ with $v_j(y)=T^{w_ij}u_j(y)$ for every $y\in A$ and $j<k$, then
there is $h\in\Gc_{i-1}^{w,D}$ with $\agree_\mu(u,h)\ge\nu(A)$.
\end{lemma}

\begin{proof}\begin{enumerate}
\item Since $\nu(A)>\rho$ and $\nu\le1$ give $|A|\ge\nu(A)|L_i|>\rho|L_i|=k^{m-i}$
(\ref{def:domains}), \ref{lem:div} with $\gamma=w_ij$, $v=v_j$ and this $A$ gives $v_j=T^{w_ij}p_j$
with $p_j\in\Gc_i^{w,D-w_ij}$; cancelling in $\F[T]$ gives $u_j=p_j$ on $A$.
\item By \ref{lem:unfold}, $h:=\sum_{j<k}x^jp_j(x^k)$ lies in $\Gc_{i-1}^{w,D}$ with components
$p_j$.
\item By \ref{lem:fold-invert}, $u(x)=\sum_jx^ju_j(x^k)=\sum_jx^jp_j(x^k)=h(x)$ for every
$x\in\fib(y)$, $y\in A$.
\item By \ref{lem:lift}, $\agree_\mu(u,h)\ge\nu(A)$.
\end{enumerate}\end{proof}

\begin{lemma}[Round lemma]\label{lem:round}\lean{SymFRI.round}\uses{hyp:ca, cor:share, lem:round-up, lem:ext, lem:ext-agree, lem:gran}
Assume $\mathrm H(\alpha_*,\Theta)$ and $\alpha_*>\rho$. Let $1\le i\le m$, let $u:L_{i-1}\to\F[T]_{\le d_{i-1}}$,
and let $\mu^{(i-1)}$ be a weight on $L_{i-1}$ with values in $\frac1{k^{i-1}}\Z$; let
$\nu^{(i)}$ be its fibre average (\ref{def:fibre-avg}). Put
\[
\beta_i:=\max\Big(\agree_{\mu^{(i-1)}}\big(u,\Gc_{i-1}^{w,D}\big),\ \alpha_*\Big),\qquad
S_i:=\Big\{r\in\F:\ \agree_{\nu^{(i)}}\big(\Fold_{rT^{w_i}}(u),\Gc_i^{w,D}\big)>\beta_i\Big\}.
\]
Then $|S_i|\le\Theta\big(|L_i|,k^i,k-1\big)$.
\end{lemma}

\begin{proof}\begin{enumerate}
\item Assume $|S_i|>\Theta(|L_i|,k^i,k-1)$, let $u_j$ be the components of $u$ and put $\tilde u_j:=T^{w_ij}u_j$; by \ref{def:fold}, $\Fold_{rT^{w_i}}(u)=\sum_{j<k}r^j\tilde u_j$ for every $r\in\F$.
\item By \ref{lem:gran}, $S_i=\{r:\agree_{\nu^{(i)}}(\sum_jr^j\tilde u_j,\Gc_i^{w,D})\ge\beta_i^+\}$, and $\beta_i^+\ge\alpha_*$.
\item By \ref{lem:ext} with $[\E:\F]:=1+\max(D,d_{i-1}+w_i(k-1))$, the $\tilde u_j$ become words over $\E$ and $\Gc_i^{w,D}$ an $\F$-linear subspace $W$ of $\RS[\E,L_i,k^{m-i}]$, of rate $\rho$; by \ref{lem:ext-agree} the agreement of $\varphi(\tilde u_j)$ with $W$ is that of $\tilde u_j$ with $\Gc_i^{w,D}$.
\item So \ref{cor:share} with $n=|L_i|$, $M=k^i$, $l=k-1$ gives $v_0,\dots,v_{k-1}\in\Gc_i^{w,D}$ and $A\subseteq L_i$ with $\nu^{(i)}(A)\ge\beta_i^+$ and $\tilde u_j=v_j$ on $A$.
\item Since $\beta_i^+\ge\alpha_*>\rho$, \ref{lem:round-up} with $\mu:=\mu^{(i-1)}$, $\nu:=\nu^{(i)}$ and this $A$ gives $h\in\Gc_{i-1}^{w,D}$ with $\agree_{\mu^{(i-1)}}(u,h)\ge\nu^{(i)}(A)\ge\beta_i^+>\beta_i$, contradicting the definition of $\beta_i$ as a maximum over $\Gc_{i-1}^{w,D}$.
\end{enumerate}\end{proof}

\section{Soundness}

\begin{lemma}[Good challenges keep the agreement]\label{lem:sound-fold}\lean{SymFRI.sound_fold}\uses{hyp:ca, lem:start, lem:round, lem:no-increase}
Assume $\mathrm H(\alpha_*,\Theta)$ with $\alpha_*>\rho$, let $\alpha\ge\alpha_*$, and let
$\agree(u_0,C_{w,D})\le\alpha$. Fix a prover and challenges $r_1,\dots,r_m\in\F$, let
$u_1,\dots,u_m$ be the prover's words, and let $S_i\subseteq\F$ be the set of \ref{lem:round}
for $u:=u_{i-1}$ and $\mu^{(i-1)}$; it depends only on $r_1,\dots,r_{i-1}$. If $r_i\notin S_i$
for every $1\le i\le m$, then $\agree_{\mu^{(m)}}(u_m,\Gc_m^{w,D})\le\alpha$.
\end{lemma}

\begin{proof}
By \ref{lem:start}, $\agree_{\mu^{(0)}}(u_0,\Gc_0^{w,D})=\agree(u_0,C_{w,D})\le\alpha$. For
$1\le i\le m$, $r_i\notin S_i$ means by \ref{lem:round} that
$\agree_{\nu^{(i)}}(\Fold_{c_i}(u_{i-1}),\Gc_i^{w,D})\le\beta_i$, and $\beta_i$ is the maximum of
$\agree_{\mu^{(i-1)}}(u_{i-1},\Gc_{i-1}^{w,D})$ and $\alpha_*$; with \ref{lem:no-increase} this
gives
$\agree_{\mu^{(i)}}(u_i,\Gc_i^{w,D})\le\max(\agree_{\mu^{(i-1)}}(u_{i-1},\Gc_{i-1}^{w,D}),\alpha_*)$,
and induction from the base case, using $\alpha\ge\alpha_*$, gives the claim.
\end{proof}

\begin{theorem}[Soundness, reduction form]\label{thm:reduction}\lean{SymFRI.reduction}\uses{lem:sound-fold, lem:bad-count, ext:claim85}
Assume $\mathrm H(\alpha_*,\Theta)$ with $\alpha_*>\rho$, let $\alpha\ge\alpha_*$, and suppose
$\agree(u_0,C_{w,D})\le\alpha$. Then for every prover the test $\Pi_{w,D}$ accepts with
probability at most
\[
\frac1{|\F|}\sum_{i=1}^{m}\Theta\big(|L_i|,k^i,k-1\big)+\alpha^s .
\]
\end{theorem}

\begin{proof}
The prover's words are fixed; the challenges $(r_1,\dots,r_m)$ are uniform on $\F^m$. By
\ref{lem:bad-count} with $S_i$ as in \ref{lem:round} and $\Theta_i:=\Theta(|L_i|,k^i,k-1)$, the
tuples with $r_i\in S_i$ for some $i$ number at most $|\F|^{m-1}\sum_i\Theta_i$, so they have
probability at most $\frac1{|\F|}\sum_i\Theta_i$. Outside them \ref{lem:sound-fold} gives
$\agree_{\mu^{(m)}}(u_m,\Gc_m^{w,D})\le\alpha$; since $u_m=g\in\Gc_m^{w,D}$ and $\Gc_m^{w,D}$ is the
set of constant words with value in $\F[T]_{\le D}$, \ref{ext:claim85} bounds the probability that
one query accepts by $\alpha$ and the $s$ repetitions by $\alpha^s$. Adding the two bounds gives the
claim.
\end{proof}

\begin{consequence}[Soundness at the Johnson radius]\label{cor:johnson}\lean{SymFRI.johnson}\uses{prop:H1, thm:reduction}
Let $\kappa\ge3$, $\alpha\ge\alpha_0:=\sqrt\rho\,(1+\frac1{2\kappa})$, and
suppose $\agree(u_0,C_{w,D})\le\alpha$. Then the test accepts with probability at most
$\varepsilon_C+\alpha^s$, where
\[
\varepsilon_C:=\frac1{|\F|}\Big(\frac{(1+\frac1{2\kappa})^7\kappa^7\,N^2}{3(k+1)\rho^{3/2}}
+m(k-1)\frac{(2\kappa+1)(N+1)}{\sqrt\rho}\Big).
\]
\end{consequence}

\begin{proof}
Use \ref{prop:H1} in \ref{thm:reduction}: for each $i$ the quantity
$\Theta_J(|L_i|,k^i,k-1)$ equals $(k-1)\max(a_J|L_i|^2,b_J(N+1))$, with $a_J$ and $b_J$ the two
constants of $\Theta_J$. Here $M=k^i$, so $Mn=N$. Hence
\[
\sum_{i=1}^{m}\Theta_J(|L_i|,k^i,k-1)\ \le\ a_J\sum_{i=1}^{m}(k-1)|L_i|^2+b_J\,m(k-1)(N+1),
\]
and $\sum_i(k-1)|L_i|^2=(k-1)N^2\sum_ik^{-2i}=\frac{N^2(1-k^{-2m})}{k+1}\le\frac{N^2}{k+1}$.
\end{proof}

\begin{consequence}[Soundness in the unique-decoding regime]\label{cor:ud}\lean{SymFRI.unique_decoding}\uses{prop:H2, thm:reduction}
Let $\alpha>\frac{1+\rho}2$ and suppose $\agree(u_0,C_{w,D})\le\alpha$. Then the test accepts
with probability at most $\varepsilon_{\mathrm U}+\alpha^s$, where
\[
\varepsilon_{\mathrm U}:=\frac{m(k-1)(N+1)}{|\F|}\Big(\frac1{\eta_{\mathrm U}}+\frac3{\sqrt\rho}\Big),
\qquad \eta_{\mathrm U}=\min\Big(\frac{(1-\sqrt\rho)^2}2,\frac{\sqrt\rho}{20}\Big).
\]
\end{consequence}

\begin{proof}
Use \ref{prop:H2} with $\alpha_*=\alpha$ in \ref{thm:reduction}: each of the $m$ terms
$\Theta_U(|L_i|,k^i,k-1)$ equals $(k-1)(N+1)(\frac1{\eta_{\mathrm U}}+\frac3{\sqrt\rho})$ by \ref{def:domains}.
\end{proof}

\begin{remark}[Individual degrees]
Since $C_{w,D}\subseteq\RS[\F,L_0,k^m]$, acceptance also certifies $\deg_{x_l}<k$ for all $l$;
taking $w=e_l$ and $D=b<k$ certifies $\deg_{x_l}f\le b$.
\end{remark}

\begin{remark}[Characteristic]\label{rem:char}\notready\uses{ext:bciks72}
\ref{prop:H1} inherits from \ref{ext:bciks72} the separability convention of
[BCIKS20, \S5.2.3] (the characteristic of $\F$ exceeds $D_Y$, the bound on the $Y$-degree of the interpolating polynomial $Q(X,Y,Z)$ of [BCIKS20, \S5], which satisfies $D_Y<\frac{2\kappa+1}{2\sqrt\rho}$); the
general case, with inseparable factors, is treated in [BCIKS20, Appendix~C].
\end{remark}

\section{Completeness}

{\renewcommand{\theproposition}{P1}
\begin{proposition}[Completeness and the terminal polynomial]\label{prop:complete}\lean{SymFRI.complete}\uses{def:protocol, lem:stab-cap, cor:terminal, cor:box}
Let $f\in\F[x_1,\dots,x_m]_{<k}$ with $\deg_wf\le D$, and let the prover be honest.
\begin{enumerate}
\item[(a)] Every format condition and every query check of $\Pi_{w,D}$ holds for every
choice of challenges; the verifier accepts with probability $1$.
\item[(b)] The final message is $g(T)=f(r_1T^{w_1},\dots,r_mT^{w_m})$.
\item[(c)] if $f\ne0$, then $\deg g\le\deg_wf$, and
$\Pr_r\big[\deg g<\deg_wf\big]\le\frac{m(k-1)}{|\F|}$.
\end{enumerate}
\end{proposition}}

\begin{proof}
(a) $u_0=U_f|_{L_0}\in C_{w,D}$. By \ref{lem:stab-cap}, each honest $u_i$ takes values of
$T$-degree $\le d_i$, so the formats hold, and each check holds because
$u_i=\Fold_{c_i}(u_{i-1})$ by definition.
(b) is \ref{cor:terminal}.
(c) The coefficient of $T^{\deg_wf}$ in $g$ is $f^{\mathrm{top}}(r)$, where
$f^{\mathrm{top}}:=\sum_{\langle w,\alpha\rangle=\deg_wf}f_\alpha x^\alpha\ne0$. It satisfies
$\deg_{x_l}f^{\mathrm{top}}<k$ for every $l$, so it vanishes with probability $\le m(k-1)/|\F|$ by
\ref{cor:box}, and no larger power of $T$ occurs.
\end{proof}

\section{Several degree constraints}

Let $p\ge1$, let $w^{(1)},\dots,w^{(p)}\in\N^m$ be weights and $D_1,\dots,D_p$ bounds, and
write $\mathbf w:=(w^{(1)},\dots,w^{(p)})$, $\mathbf D:=(D_1,\dots,D_p)$ and $d^{(\iota)}_i:=\min\big(D_{\iota},(k-1)\sum_{l\le i}w^{(\iota)}_l\big)$. The target code is $C_{\mathbf w,\mathbf D}:=\bigcap_{\iota}C_{w^{(\iota)},D_{\iota}}$.

\begin{remark}[Several constraints in one run]\label{rem:multi}\lean{SymFRI.multi}\uses{lem:inj, lem:fold-sub, lem:stab, lem:div, lem:unfold, lem:round, thm:reduction}
Replace $\F[T]$ by $\F[T_1,\dots,T_p]$, the challenge $r_iT^{w_i}$ by
$r_i\prod_{\iota}T_{\iota}^{\,w^{(\iota)}_i}$, the graded degree by the vector of degrees
$\big(a_{\iota}+\sum_{l>i}w^{(\iota)}_l\alpha_l\big)_{\iota}$ of a monomial $\prod_{\iota}T_{\iota}^{a_{\iota}}x^\alpha$, the
condition $\gdeg\le D$ by ``$\le D_{\iota}$ in every coordinate'', the formats $\F[T]_{\le d_i}$ by
$\prod_{\iota}\F[T_{\iota}]_{\le d^{(\iota)}_i}$, and the final check $\deg g\le D$ by $\deg_{T_{\iota}}g\le D_{\iota}$
for all $\iota$. Then \ref{lem:inj}, \ref{lem:fold-sub}, \ref{lem:stab}, \ref{lem:div},
\ref{lem:unfold}, \ref{lem:round} and \ref{thm:reduction} extend verbatim to
$C_{\mathbf w,\mathbf D}$, with $[\E:\F]:=\prod_{\iota}\big(1+\max(D_{\iota},\,d^{(\iota)}_{i-1}+w^{(\iota)}_i(k-1))\big)$
in item~3 of \ref{lem:round}.


The proofs use only three properties of $\F[T]$: it is an integral domain (cancellation in
\ref{lem:unfold} and \ref{lem:round}); divisibility by a monomial is read off coefficientwise
(\ref{lem:div} and \ref{lem:round}, item~4); and the value spaces are finite-dimensional over
$\F$, which is all \ref{lem:ext} needs. All three hold for $\F[T_1,\dots,T_p]$ with bounded
degrees, with the bound $t$ above.
\end{remark}

\begin{remark}[Separate runs on one commitment]\label{rem:separate}\lean{SymFRI.separate}\uses{rem:multi}
Running $\Pi_{w^{(\iota)},D_{\iota}}$ separately on the same $u_0$ costs the sum of the $p$ runs instead
of the product $t$. Below the unique-decoding radius it certifies the same statement: if
$\agree(u_0,C_{w^{(\iota)},D_{\iota}})>\frac{1+\rho}2$ for all $\iota$, the codewords achieving these
agreements are all the unique closest codeword of $\RS[\F,L_0,k^m]$, which therefore lies in
$C_{\mathbf w,\mathbf D}$. Above that radius, separate runs only certify closeness to each
$C_{w^{(\iota)},D_{\iota}}$, possibly through different codewords of the list; \ref{rem:multi} is then
the right tool.
\end{remark}

\section{The bivariate setting of {FRISS}}

\begin{definition}[Bivariate setting]\label{def:avss}\lean{SymFRI.Vt}\leanok\uses{def:domains, def:kron, def:codes}
Let $\tau,t\in\N$ with $2^\tau\ge t+1$. Set $k:=2$ and $m:=2\tau$, and let $N,L_0,\dots,L_m$ be
as in \ref{def:domains}, so $2^{2\tau}=k^m\mid N$, and $\rho=k^m/N=2^{2\tau}/N$ by
\ref{def:domains}. Here $X,Y$ are two
indeterminates, $\deg_X$ and $\deg_Y$ the degrees in them, and $C_{w^X,t}$ is the degree code of
\ref{def:codes} for the weight $w^X$ below and $D=t$. For
$f\in\F[X,Y]$ with $\deg_Xf,\deg_Yf<2^\tau$, let $f'\in\F[x_1,\dots,x_{2\tau}]_{<2}$ be the
multilinear polynomial with
\[
f(X,Y)=f'\big(X,X^2,\dots,X^{2^{\tau-1}},\,Y,Y^2,\dots,Y^{2^{\tau-1}}\big),
\]
and put $\hat f:=U_{f'}$; substituting $x_l=Z^{2^{l-1}}$ in $U_{f'}$ gives
$\hat f(Z)=f(Z,Z^{2^\tau})$. The map $(\cdot)'$ is $\F$-linear, so $(f+\theta e)'=f'+\theta e'$ and
$\hat h=\hat f+\theta\hat e$ for $h:=f+\theta e$. The weights are
\[
w^X:=(1,2,\dots,2^{\tau-1},0,\dots,0),\qquad
w^Y:=(0,\dots,0,1,2,\dots,2^{\tau-1})\in\N^{2\tau},
\]
the box of the intermediate degree bound is $V_t:=\{f\in\F[X,Y]:\ \deg_Xf\le t,\ \deg_Yf\le t\}$,
$u_f,u_e:L_0\to\F$ are the committed words, honestly $u_f=\hat f|_{L_0}$ and
$u_e=\hat e|_{L_0}$, and $\delta_0\in[0,1]$ bounds their distance to these words.
\end{definition}

\begin{lemma}[Binary digits: individual degree is weighted degree]\label{lem:avss-bits}\lean{SymFRI.avss_bits}\uses{def:avss, def:wdeg}
Let $f\in\F[X,Y]$ with $\deg_Xf,\deg_Yf<2^\tau$, and let $f',w^X,w^Y$ be as in
\ref{def:avss}. Then
\[
\deg_Xf=\deg_{w^X}f',\qquad \deg_Yf=\deg_{w^Y}f'.
\]
\end{lemma}

\begin{proof}\begin{enumerate}
\item Under the substitution of \ref{def:avss} the monomial $x^\alpha$,
$\alpha\in\{0,1\}^{2\tau}$, maps to $X^{a(\alpha)}Y^{b(\alpha)}$ with
$a(\alpha):=\sum_{l\le\tau}\alpha_l2^{l-1}=\langle w^X,\alpha\rangle$ and
$b(\alpha):=\sum_{l>\tau}\alpha_l2^{l-\tau-1}=\langle w^Y,\alpha\rangle$, and its coefficient
is unchanged.
\item By uniqueness of the binary expansion, $\alpha\mapsto(a(\alpha),b(\alpha))$ is a bijection
from $\{0,1\}^{2\tau}$ to $\{0,\dots,2^\tau-1\}^2$; hence $f'_\alpha=f_{a(\alpha),b(\alpha)}$ for every
$\alpha$, and every pair $(a,b)$ with $f_{a,b}\ne0$ is of the form $(a(\alpha),b(\alpha))$.
\item Therefore $\deg_Xf=\max\{\langle w^X,\alpha\rangle:\ f'_\alpha\ne0\}=\deg_{w^X}f'$, and
likewise for $Y$.
\end{enumerate}\end{proof}

\begin{remark}[The $Y$ direction]\label{rem:avss-y}\notready\uses{def:avss, def:kron, lem:avss-bits}
With $X$ the low variable, $\deg_Yf\le t$ if and only if $\deg\hat f<2^\tau(t+1)$. Indeed
$\deg_Yf=\deg_{w^Y}f'$ by \ref{lem:avss-bits}, and by \ref{def:kron}
$\hat f(Z)=f(Z,Z^{2^\tau})=\sum_{a,b}f_{a,b}Z^{a+2^\tau b}$ with $0\le a<2^\tau$; hence
$\deg\hat f=\max\{a+2^\tau b:\ f_{a,b}\ne0\}$, and for $0\le a<2^\tau$ one has
$b\le t\iff a+2^\tau b<2^\tau(t+1)$.
\end{remark}

\section{The masked degree test}

\begin{definition}[Masked degree test]\label{con:avss-mask}\uses{def:avss, def:protocol}
Let $\tau,t$ and $f\in\F[X,Y]$ with $\deg_Xf,\deg_Yf<2^\tau$ be as in \ref{def:avss}. The
dealer samples $e$ uniformly from $V_t$ and commits to $\hat e$ with salted Merkle leaves: the
leaf of $x$ hashes $x$, $\hat e(x)$ and an independent uniform salt. The verifier sends
$\theta\in\F$. Put $h:=f+\theta e$ and $u_h:=u_f+\theta u_e$. The two parties run
$\Pi_{w^X,t}$ (\ref{def:protocol}) with $k=2$, $m=2\tau$, $D=t$ and $u_0:=u_h$, the verifier
reading $u_h(x)$ by opening $u_f(x)$ and $u_e(x)$; the $s$
repetitions use as query points $\xi_0$ the $s$ query positions at which the commitment already opens $\hat f$, in order, so that $s$ equals the number of queries of the commitment and each $\xi_0$ is uniform on $L_0$.
\end{definition}

\begin{remark}[The threshold is not enforced]\label{rem:avss-threshold}\notready\uses{def:avss}
Let $2^\tau\ge t+2$ and let $f\in\F[X,Y]$ satisfy $\deg_Xf,\deg_Yf<2^\tau$ and
$\deg_Xf(X,0)>t$; for instance $f=X^{t+1}$ if $t+1\le2^\tau-1$. PolyFRIM's binding certifies
only $\deg_Xf,\deg_Yf<2^\tau$ {[PolyFRIM, Def.~2.3, App.~B, App.~D]}, and the dealer's messages
and proofs are those of an honest dealer for this $f$, so every check of FRISS
{[PolyFRIM, Alg.~2--4]} passes. The Rec\_strong value of party $i$ is $\sigma_i:=f(\omega_i,0)$, at
distinct points $\omega_i\in\F$ {[PolyFRIM, Alg.~3--4]}. For a set $A$ of $t+1$ parties let
$P_A\in\F[X]$ be the interpolant of $\{(\omega_i,\sigma_i):\ i\in A\}$, of degree at most $t$, and put
$\sigma:=f(z,0)$. Then one of two cases occurs.
\begin{enumerate}
\item If $P_A\ne P_B$ for two sets $A,B$ of $t+1$ parties, then $P_A,P_B$ are distinct of
degree $\le t$, so by \ref{lem:roots} with $D=t+1$ one has $P_A(z)\ne P_B(z)$ outside at most $t$
values of $z$: the parties may output different secrets.
\item If $P_A=P$ for every $A$, then $P$ agrees with $f(\cdot,0)$ on every $\omega_i$ while
$P\ne f(\cdot,0)$, and both have degree $\le2^\tau-1$, so by \ref{lem:roots} with $D=2^\tau$ one has
$P(z)\ne\sigma$ outside at most $2^\tau-1$ values of $z$: every party outputs a secret other than
$s$.
\end{enumerate}
The same argument applies with $X$ and $Y$ exchanged, to the shares $\hat f_i=f(\omega_i,Y)$
interpolated from $t+1$ valid Echo values {[PolyFRIM, Alg.~3, step 16]}.

This does not contradict {[PolyFRIM, Thm.~4.1]}: Definition~2.4 there constrains the
reconstructed value only when the dealer is honest (Correctness), and its Agreement property
requires that all honest parties terminate, not that they output the same value. The remark
shows that the stronger requirement, that a malicious dealer be bound to a single secret, is not
met without a degree test.
\end{remark}


\begin{lemma}[One bad challenge]\label{lem:avss-bad}\lean{SymFRI.avss_bad}
Let $f,e\in\F[X,Y]$ with $\deg_Xf>t$. Then $\deg_X(f+\theta e)\le t$ for at most one
$\theta\in\F$.
\end{lemma}

\begin{proof}
Let $\theta\ne\theta'$ and suppose $\deg_X(f+\theta e),\deg_X(f+\theta'e)\le t$. Then
$\deg_X\big((\theta-\theta')e\big)\le t$, and $\theta-\theta'\ne0$ is a scalar, so $\deg_Xe\le t$;
hence $\deg_Xf=\deg_X\big((f+\theta e)-\theta e\big)\le t$, contradicting $\deg_Xf>t$.
\end{proof}

\section{Distance of the masked word}

\begin{lemma}[The commitments add up]\label{lem:avss-close}\lean{SymFRI.avss_close}\uses{def:avss, def:agree}
Let $f\in\F[X,Y]$ with $\deg_Xf,\deg_Yf<2^\tau$, let $e\in\F[X,Y]$ with $\deg_Xe,\deg_Ye<2^\tau$,
let $\theta\in\F$ and $u_f,u_e:L_0\to\F$, and put $h:=f+\theta e$, $u_h:=u_f+\theta u_e$ and $\hat h:=U_{h'}$ with $h'$ as in
\ref{def:avss}. If $\Delta(u_f,\hat f|_{L_0})\le\delta_0$ and
$\Delta(u_e,\hat e|_{L_0})\le\delta_0$, then $\agree(u_h,\hat h|_{L_0})\ge1-2\delta_0$.
\end{lemma}

\begin{proof}
Outside the union of the two disagreement sets, of size at most $2\delta_0N$, one has
$u_h(x)=u_f(x)+\theta u_e(x)=\hat f(x)+\theta\hat e(x)=\hat h(x)$.
\end{proof}

\begin{lemma}[Distance to the degree code]\label{lem:avss-dist}\lean{SymFRI.avss_dist}\uses{def:avss, con:avss-mask, lem:avss-bits, lem:avss-close, lem:inj, lem:roots, def:codes, def:agree}
Let $f\in\F[X,Y]$ with $\deg_Xf,\deg_Yf<2^\tau$, let $e\in\F[X,Y]$ with $\deg_Xe,\deg_Ye<2^\tau$, let $\theta\in\F$, put
$h:=f+\theta e$ and $u_h:=u_f+\theta u_e$, and assume
$\Delta(u_f,\hat f|_{L_0})\le\delta_0$ and $\Delta(u_e,\hat e|_{L_0})\le\delta_0$. If
$\deg_Xh>t$, then
\[
\agree(u_h,C_{w^X,t})\le\rho+2\delta_0 .
\]
\end{lemma}

\begin{proof}\begin{enumerate}
\item \ref{lem:avss-close} gives $\agree(u_h,\hat h|_{L_0})\ge1-2\delta_0$.
\item Let $c=U_q|_{L_0}\in C_{w^X,t}$, so $q\in\F[x_1,\dots,x_{2\tau}]_{<2}$ and
$\deg_{w^X}q\le t<\deg_Xh=\deg_{w^X}h'$ by \ref{lem:avss-bits}; then $q\ne h'$, and
\ref{lem:inj} applied to the $T$-free polynomials $q$ and $h'$ gives
$\hat q=U_q\ne U_{h'}=\hat h$.
\item Two distinct polynomials of degree $<2^{2\tau}$ agree on fewer than $2^{2\tau}=\rho N$
points of $L_0$ (\ref{lem:roots} with $D=2^{2\tau}$), so
$\agree(\hat h|_{L_0},c)<\rho$, that is $\Delta(\hat h|_{L_0},c)>1-\rho$.
\item By \ref{def:agree}, $\Delta(u_h,c)\ge\Delta(\hat h|_{L_0},c)-\Delta(u_h,\hat h|_{L_0})
>1-\rho-2\delta_0$ by steps 1 and 3, so $\agree(u_h,c)<\rho+2\delta_0$; taking the maximum over
$c\in C_{w^X,t}$ gives $\agree(u_h,C_{w^X,t})\le\rho+2\delta_0$.
\end{enumerate}\end{proof}

\section{Soundness of the masked degree test}

{\renewcommand{\theproposition}{P2}
\begin{proposition}[Soundness]\label{prop:avss-sound}\lean{SymFRI.avss_sound}\uses{def:avss, con:avss-mask, lem:avss-bad, lem:avss-dist, cor:ud}
Let $\tau,t,u_f,u_e,\delta_0$ be as in \ref{def:avss}, let $\delta_0\le(1-\rho)/4$, let
$f,e\in\F[X,Y]$ have $\deg_Xf,\deg_Xe,\deg_Yf,\deg_Ye<2^\tau$, where $e$ is the mask committed by the
dealer (not necessarily in $V_t$), and let $\deg_Xf>t$. Then for every $\alpha_*>(1+\rho)/2$ the masked test
\ref{con:avss-mask}, run with $w=w^X$, $k=2$, $m=2\tau$ and $D=t$, accepts with probability at
most
\[
\frac1{|\F|}+\varepsilon_{\mathrm U}+\alpha_*^s,
\]
where $\varepsilon_{\mathrm U}$ is the error of \ref{cor:ud} for these parameters, that is
$\varepsilon_{\mathrm U}=2\tau\,(N+1)\big(\eta_{\mathrm U}^{-1}+3/\sqrt\rho\big)/|\F|$.
\end{proposition}}

\begin{proof}\begin{enumerate}
\item By \ref{lem:avss-bad} the set $B$ of $\theta\in\F$ with $\deg_X(f+\theta e)\le t$ has at
most one element, so $\Pr[B]\le1/|\F|$.
\item Outside $B$ one has $\deg_Xh>t$, and \ref{lem:avss-dist} gives
$\agree(u_h,C_{w^X,t})\le\rho+2\delta_0\le\rho+\frac{1-\rho}2=\frac{1+\rho}2<\alpha_*$.
\item On that event the transcript is a run of $\Pi_{w^X,t}$ with $u_0=u_h$ fixed by the two
commitments and with $\agree(u_0,C_{w^X,t})\le\alpha_*$, so \ref{cor:ud} bounds the acceptance
probability by $\varepsilon_{\mathrm U}+\alpha_*^s$.
\item Adding the bound of step 1 gives $\frac1{|\F|}+\varepsilon_{\mathrm U}+\alpha_*^s$.
\end{enumerate}\end{proof}

\section{Secrecy of the masked degree test}

\begin{lemma}[The masked value is uniform]\label{lem:avss-uniform}\lean{SymFRI.avss_uniform}\uses{def:avss}
Let $f,e\in V_t$ with $e$ uniform and independent of $f$ and of $\theta\in\F^*$. Then
$h:=f+\theta e$ is uniform on $V_t$ and independent of $f$ and of $\theta$.
\end{lemma}

\begin{proof}
The map $e\mapsto f+\theta e$ is a translation composed with a scaling by $\theta\ne0$, hence a
bijection of $V_t$; a uniform variable pushed through a bijection stays uniform, and as a function
of $e$ alone $h$ is independent of $f$ and of $\theta$.
\end{proof}

\begin{lemma}[Secrecy]\label{lem:avss-secret}\notready\uses{def:avss, con:avss-mask, def:protocol, lem:avss-uniform}
Let $f\in V_t$ and $\theta\in\F^{*}$, and let $e$ be uniform on $V_t$, independent of $f$, of
$\theta$ and of the randomness of the PCS. Then
\begin{enumerate}
\item[(a)] every message of the masked test \ref{con:avss-mask} is a function of $h=f+\theta e$,
$\theta$, the challenges $r_1,\dots,r_m$ and the query positions;
\item[(b)] given $\theta$, the challenges, the positions and the values of $\hat f$ at those
positions, a simulator produces a transcript distributed exactly as the honest one.
\end{enumerate}
\end{lemma}

\begin{proof}
(a) The messages of the dealer are $u_h$, its folds on the queried fibres and $g$, which
\ref{def:protocol} determines by $h$ and $r_1,\dots,r_m$; the openings at a queried position $x$ are
$\hat h(x)=\hat f(x)+\theta\hat e(x)$, so opening $\hat e$ reveals
$\hat e(x)=\big(\hat h(x)-\hat f(x)\big)/\theta$ and requires $\hat f$ only at the positions --- the
one use of $\theta\ne0$ besides \ref{lem:avss-uniform}.

(b) Sample $h$ uniform on $V_t$, which is legitimate by \ref{lem:avss-uniform}; output the
transcript of (a) and the $\hat e$-openings; the leaves that are not opened are hidden by the salts
of \ref{con:avss-mask}, and a simulator holding the values of $\hat f$ at the positions produces the
same distribution. The case $\theta=0$, the only one in which \ref{lem:avss-uniform} fails, has
probability $1/|\F|$ in \ref{con:avss-mask}.
\end{proof}

\begin{remark}[Leaked evaluations]\label{rem:avss-leak}\notready\uses{con:avss-mask}
If the query points of the masked test cannot be chosen among the positions of the commitment,
then each of the $s$ repetitions opens $u_f$ at the $k=2$ points of a fibre, so the test leaks
$2s$ further evaluations of $\hat f$, and the threshold of {[PolyFRIM, \S4, App.~D]} becomes
$t=c+4s$.
\end{remark}

\section{Cost of the masked degree test}

\begin{lemma}[The caps for $w^X$]\label{lem:avss-caps}\uses{def:avss, def:protocol}
Let $k=2$, $w=w^X$, $D=t$ and $t\le2^\tau-1$, and let
$d_i=\min\big(D,(k-1)\sum_{l\le i}w_l\big)$. Then $d_i=\min(t,2^i-1)$ for $0\le i\le\tau$ and
$d_i=t$ for $\tau<i\le2\tau$.
\end{lemma}

\begin{proof}
For $i\le\tau$, $(k-1)\sum_{l\le i}w_l=\sum_{l\le i}2^{l-1}=2^i-1$; for $i>\tau$ the sum is
$2^\tau-1$, so $t\le2^\tau-1$ gives $d_i=\min(t,2^\tau-1)=t$.
\end{proof}

{\renewcommand{\theproposition}{P3}
\begin{proposition}[Cost of one repetition]\label{prop:avss-cost}\lean{SymFRI.avss_cost}\uses{def:avss, con:avss-mask, def:protocol, lem:avss-caps}
Let $w=w^X$, $D=t$, $k=2$, $m=2\tau$, and let $d_i=\min\big(D,(k-1)\sum_{l\le i}w_l\big)$ be
the caps of \ref{def:protocol}. Then $d_i=\min(t,2^i-1)$ for $i\le\tau$ and $d_i=t$ for
$i>\tau$, and one repetition of the masked test \ref{con:avss-mask} reads
\[
2\sum_{i=1}^{2\tau}(d_{i-1}+1)\le2\big(2^\tau-1+\tau(t+1)\big)<2\big(2^\tau+\tau(t+1)\big)
\]
field elements. With $\tau=\lceil\log_2(t+1)\rceil$ this is at most
$2(t+1)\big(3+\log_2(t+1)\big)$, hence $O(t\log t)$; over $s$ repetitions the total is
$s\cdot O(t\log t)$, which for fixed $s$ and $t=\Theta(n)$ is $O(n\log n)$, inside the
per-party cost $O(n\log^2n)$ of {[PolyFRIM, \S4]}.
\end{proposition}}

\begin{proof}\begin{enumerate}
\item By \ref{lem:avss-caps}, $d_j+1=\min(t+1,2^j)\le2^j$ for $j<\tau$, and $d_j=t$ for
$j\ge\tau$.
\item Splitting the sum at $\tau$,
$\sum_{i=1}^{2\tau}(d_{i-1}+1)\le\sum_{j=0}^{\tau-1}2^j+\sum_{j=\tau}^{2\tau-1}(t+1)=2^\tau-1+\tau(t+1)$, with equality if $2^{\tau-1}\le t+1$.
\item The query cost of the protocol of \ref{def:protocol} is $k\sum_{i=1}^m(d_{i-1}+1)$, which is
twice the sum of step 2, giving the stated inequalities.
\item If $\tau=\lceil\log_2(t+1)\rceil$, then $2^\tau<2(t+1)$ and $\tau\le1+\log_2(t+1)$, so the
count is less than $2(t+1)\big(2+1+\log_2(t+1)\big)$.
\item Multiplying by $s$ and using $t=\Theta(n)$ and $s=O(1)$ gives $O(n\log n)\subseteq
O(n\log^2n)$.
\end{enumerate}\end{proof}

\section{Total degree: lines and applications}

\begin{lemma}[Restriction to a line]\label{lem:line-total}\lean{SymFRI.line_total}\uses{def:wdeg, lem:subst}
Let $m\ge1$, $f\in\F[x_1,\dots,x_m]$ and $\mathbf a,\mathbf b\in\F^m$. Then
\[
f(\mathbf a+T\mathbf b):=f(a_1+Tb_1,\dots,a_m+Tb_m)\in\F[T]
\]
has degree at most $\deg_{\mathbf 1}f$.
\end{lemma}

\begin{proof}
\ref{lem:subst} with $R=\F$, $w=\mathbf 1$, $I=\{1,\dots,m\}$ and $g_l=a_l+Tb_l$, of degree
$\le1=w_l$, gives $\gdeg_{\mathbf 1}f(\mathbf a+T\mathbf b)\le\deg_{\mathbf 1}f$; a polynomial of
$\F[T]$ has no other variables, so its graded degree is its $T$-degree.
\end{proof}

\begin{remark}[Two evaluation claims reduced to one]\label{rem:total-app}\lean{SymFRI.total_app}\uses{lem:sz-total, lem:line-total, def:protocol}
Let $f$ with $\deg_{x_l}f<k$ for all $l$ be committed, and suppose that $\Pi_{\mathbf 1,D}$
(\ref{def:protocol}) has certified $\deg_{\mathbf 1}f\le D$. A protocol may then have to convince
the verifier of two claims $f(\mathbf a)=y_{\mathbf a}$ and $f(\mathbf b)=y_{\mathbf b}$, for
$\mathbf a,\mathbf b\in\F^m$ and $y_{\mathbf a},y_{\mathbf b}\in\F$. Let
$\ell(T):=\mathbf a+T(\mathbf b-\mathbf a)$, so that $\ell(0)=\mathbf a$ and $\ell(1)=\mathbf b$. The
prover sends a polynomial $q\in\F[T]_{\le D}$ claimed to be $f\circ\ell$; the verifier checks
$q(0)=y_{\mathbf a}$ and $q(1)=y_{\mathbf b}$, draws a uniform $T^*\in\F$ and asks for an opening of
the commitment at $\ell(T^*)\in\F^m$. If $q\ne f\circ\ell$, then $q-f\circ\ell$ is nonzero of degree
at most $D$ by \ref{lem:line-total}, so $q(T^*)\ne f(\ell(T^*))$ with probability at least
$1-D/|\F|$ by \ref{lem:sz-total} with $m=1$. The two claims are thus replaced by one, at error
$D/|\F|$ and at a cost of $D+1$ field elements. With only the individual bounds,
$\deg_{\mathbf 1}f\le m(k-1)$ is all that is known, the same step costs $m(k-1)+1$ field elements and
its error is $m(k-1)/|\F|$. At
$k=4$, $m=8$, $D=6$ the certificate gives $7$ field elements at error $6/|\F|$, against $25$
elements at $24/|\F|$ without it.
\end{remark}

\begin{remark}[Protocols whose soundness is stated in the total degree]\label{rem:total-prior}\uses{rem:total-app}
For some interactive proofs the total degree is the parameter of the soundness statement, and a
certificate of $\deg_{\mathbf 1}f\le D$ is then an input to their analysis rather than a
strengthening of it. Fold-DCS [Fold-DCS, arXiv:2504.00693] is one: "the soundness of Fold-DCS
depends on the total degree of the polynomial, not its individual degrees" (p.~2). Its error is
bounded by $\log(m)D/|\F|$ for $m$ variables and total degree $D$, against $mb/|\F|$ for the
classical sumcheck with individual degree bound $b$ (their $\mu$ and $d$) (\S3.2.2), and it constrains the commitment itself
to total degree at most $D$ (\S4.2.1). The line step of \ref{rem:total-app} is the one of the GKR
protocol [GKR, STOC 2008 / J.~ACM 62(4), 2015], which uses it to reduce the two claims of a round to
a single one [Thaler, \S4.6.4]; the subroutine, with its soundness claim, is
[Thaler, \S4.5.2, Claim~4.6].
\end{remark}
